\section[Sentence]{Sentence (Satz)} \label{sec:destce}
Sentences generally follow a Subject-Verb-Object type construction similar to English, but with a catch.  In \autoref{sec:deverb}, it was mentioned that the CV always stays in second place, and the DC is kept at the end. This rule, which will be called the second place rule (SPR) is defined as:
\begin{syntax}
    \begin{itemize}
    \item For any verb chain in any sentence, the DC always stays at the very end, and the CV always stays exactly after the first unit in the sentence
    \item If any rule requires  the CV to be shifted to the end, the CV is appended to the end of the DC
    \end{itemize}
\end{syntax}
Note that the additional rules for $V\in\mathcal{S}$  apply when appending the CV to the end of the DC. Since this  is not an aspect change, the CV will never interact with anything in the DC, even for \bt{Futur 2}.

In the context of object verbs, it must be mentioned that \bt{sein, werden} when used as main verbs do not take case $2$ nouns, but rather two case $1$ nouns:
\begin{example}
    \begin{tabular}{c|c}
        \bt{Er($1$) ist ein kriegerischer Mann($1$)} & He($1$) is a belligerent man($2$)\\\hline
        \bt{Er($1$) wird sie($1$)}& He($1$) becomes her($2$)
    \end{tabular}
\end{example}
No other verbs I know do this. As to why this is the case  is because both of them make a \bt{Perfekt} using \bt{sein}, but require objects regardless. Since \bt{sein} as a helping verb only couples with non-object verbs, the necessary object must be put in case $1$ instead to preserve the rule.
\subsection[Punctuation]{Punctuation (Zeichensetzung)}
\textgerman{\bt{Meine Kämpfe gegen die Stadt trieben mich jenseits des Landes, und in den Bereich der Verrückten. Da fand ich einen Mann, der sagte \enquote{Heil der Sonne! Wir wären nichts ohne sie!} Verwirrt? So war ich auch. Auch zu sehen auf einem Schild waren die Wörter \enquote*{Mein Leid wird bald verbrannt werden.} Plötzlich klingelte eine Glocke, und Männer, Frauen, Ältere, Jungen; selbst die Tiere des Orts liefen langsam in Richtung eines Hauses. Ich wusste, dass irgendwas ernsthaft schiefgelaufen war : Alles fühlte sich an, als wäre es tot; wie die Vernunft schon geopfert worden war.}}

This paragraph has all the punctuation used. There are no usage differences compared to English, and you will recognise the minor symbolic changes easily. Did I need to write this nonsensical paragraph to showcase this? Not at all. But being unpredictable is a mark of a good villain. The above paragraph also contains many sentence based elements which have not been introduced yet, and gives a brief outlook of how  \bt{Deutsch} normally looks.

\subsection[Conjunction]{Conjunction (Verbindungswort)}
There are two types of relational conjunctions based on how they interact with the SPR:
\begin{enumerate}
    \item \textit{Shifting}: CV appended to the end of  DC.  Can always start sentences
    \item \textit{Normal}: CV at second position. Mostly cannot start sentences
\end{enumerate}
Note that unlike verb aspect variation, shifting conjunctions do not alter the form of the CV: they only throw it to the end. For uses with cell $\vm{T}(1,2)$, the prepositional argument (if any) simply gets slammed onto the CV: \bt{abfahren/abfährt}

If any relational conjunction starts a sentence, the consequent connecting clause has its CV immediately after the comma. This occurs because the entire previous clause along with its conjunction is treated as a single  placeholder  unit, so the SPR places the next CV immediately after it:
\begin{example}
        (Dass ich hier zu diesem Zeitpunkt lebe), ist ein Wunder
\end{example}
Additional conjunctions behave exactly as they would in English, and do not effect any special changes. 
\subsubsection{Dass}
That. Can be coupled with \bt{so} to give \bt{sodass} meaning `so that'. Shifting type conjunction, and is often ignored while speaking in order to avoid exiling the CV and to improve understanding, especially when speech fillers are used:
\begin{example}
    \begin{tabular}{m{7cm}|m{7cm}}
        \bt{Ich glaube, (dass) ich mag Honig} & I believe, (that) I like honey\\\hline
        \bt{Ich unterdrücke meine Wut, sodass ich in der Gesellschaft leben kann} & I suppress my fury, so that I in \ut{the} society \ut{live can}
    \end{tabular}
\end{example}
\subsubsection{Da/Weil/Denn}
Because. Among these, \bt{da/weil} are shifting, whereas \bt{denn} is normal. \bt{Denn} is unable to start sentences, whereas the other two have no such weakness. \bt{Da} is more formal than \bt{weil}. Also, this \bt{da} obviously does not mean `there' in this context.
\begin{example}
    \begin{tabular}{m{7cm}|m{7cm}}
        \bt{Wir müssen uns beeilen, weil der Zug gleich abfährt} & We must \ut{ourselves} hurry, because the train shortly \ut{departs}\\\hline
        \bt{Ich weiß  nicht alles, denn ich bin kein Gott} & I \ut{know not} everything, because I am \ut{no} god
    \end{tabular}
\end{example}
\subsubsection{Wenn, Falls}
If/when, In case. Both of them are shifting types, and mean slightly different things. \bt{Falls} is placed at the beginning of a  conditional argument, with the other argument expressing the precaution that has been taken against that situation, which would arise, were the condition to come true:
\begin{example}
    \begin{tabular}{m{6.5cm}|m{7cm}}
        \bt{Falls das Kind aufwacht, haben wir eine Überwachungskamera im Kinderzimmer} & In case the child wakes up, \ut{have we} a security camera in the children's bedroom\\\hline
        \bt{Wir haben einen Regenschirm dabei, falls es regnet} & We have an umbrella \ut{at this point}, in case it rains
    \end{tabular}
\end{example}
\bt{Wenn} on the other hand implies cause-consequence. Although it does mean `when' sometimes, such uses are rather context dependent and not that common, since `when' is normally written using \bt{als}. In the consequent solution, \bt{wenn} exclusively means `if'. 

\bt{Wenn} has an additional construction when it starts a sentence:
\begin{itemize}
    \item Eliminate \bt{wenn} and place the CV in its place
\end{itemize}
No other conjunction allows this.
\begin{example}
    \begin{tabular}{m{6.5cm}|m{7cm}}
        \bt{Spricht man über das Wetter, vergeudet er seine Zeit}& (If) \ut{speaks} one about the weather, is wasting \ut{he} his time\\\hline
        \bt{Wasser friert, wenn seine Temperatur unter Null sinkt} & Water freezes, if its temperature under zero \ut{sinks}
    \end{tabular}
\end{example}
Adverbs can also be used as relational conjunctions, which may either be normal or shifting.  The only difference between adverbial and pure conjunctions is that adverbials always count as a unit in the SPR, whereas pure conjunctions never do. Since there are too many adverbials, only pure conjunctions will be comprehensively  listed. First, \autoref{tab:destceT1} contains all elements  $E\in\mathcal{A}_\boxdot$, where anything that counts in the SPR is marked with the exponent $\mathbb{V}$.
\begin{table}[bth]
    \centering
    \begin{subtable}{\textwidth}
        \centering
    \begin{tabular}{|c|c|}
        \hline \textit{Conjunction} & \textit{Meaning}\\\hline
        \bt{und} & and\\\hline
        \bt{oder} & or\\\hline
        \bt{beziehungsweise} & respectively / more precisely\\\hline
        \bt{aber} & but\\\hline \bt{allein} & alone\\\hline
        \bt{sondern} & but rather\\\hline
        \bt{doch}& nevertheless / however\\\hline
        \bt{als} & as (only for comparisons)\\\hline
    \end{tabular}
    \caption{Singles: $\mathcal{A}_\boxdot[\mathbb{S}]$}
\end{subtable}
\begin{subtable}{\textwidth}
    \centering
    \begin{tabular}{|c|c|m{0.8cm}|}
        \hline \textit{Conjunction} & \textit{Meaning}& \textit{Max args.}\\\hline
        \bt{$\text{entweder}^{\mathbb{V}}$\ldots oder\ldots}& either\ldots or\ldots & $\infty$\\\hline
        \bt{$\text{weder}^{\mathbb{V}}$\ldots $\text{noch}^{\mathbb{V}}$\ldots}& neither\ldots nor\ldots & $\infty$\\\hline
        \bt{nicht nur\ldots sondern\ldots auch\ldots}& not only\ldots but also\ldots (surprising) & $2$\\\hline
        \bt{sowohl\ldots als auch\ldots}& not only\ldots but also\ldots (standard) & $2$\\\hline
        \bt{$\text{zwar}^{\mathbb{V}}$\ldots aber\ldots} & indeed\ldots but\ldots (statement + context) & $2$\\\hline
        \bt{$\text{einerseits}^{\mathbb{V}}$\ldots $\text{andererseits}^{\mathbb{V}}$\ldots} & on  one hand\ldots on the other hand\ldots & $2$\\\hline
        \bt{$\text{mal}^{\mathbb{V}}$\ldots $\text{mal}^{\mathbb{V}}$\ldots}& sometimes\ldots sometimes\ldots & $\infty$\\\hline
        \bt{$\text{teils}^{\mathbb{V}}$\ldots $\text{teils}^{\mathbb{V}}$\ldots}& partly\ldots partly\ldots & $\infty$\\\hline
    \end{tabular}
    \caption{Coupled: $\mathcal{A}_\boxdot[\mathbb{O}]$}
\end{subtable}
    \caption{Additionals: $\mathcal{A}_\boxdot$}
    \label{tab:destceT1}
\end{table}

The \bt{auch} in the three part conjunction \bt{nicht nur\ldots sondern\ldots auch\ldots.} stands before the second joined unit, but \bt{sondern} always stays at the very beginning of the second clause. Hence, if there are other remaining units, they collectively form the second argument, which is otherwise the null set:
\begin{example}
    Sie bringt (nicht nur) [Wein] (sondern) gibt den Kindern (auch) [Messer]
\end{example}
Hence, there are still only two arguments, and not three. \bt{mal\ldots, teils\ldots.} can be chained. 

At least all pure elements $E\in \mathcal{R}_\boxdot$  are in \autoref{tab:destceT2}. $\mathcal{R}_\boxdot$ also contains exactly two coupled conjunctions. Any normal adverbials  are marked with the exponent $\mathbb{N}$.
\begin{table}[bth]
    \centering
    \begin{tabular}{|c|c|}
        \hline\textit{Conjunction}& \textit{Meaning}\\\hline
        \bt{als} & when\\\hline \bt{ob} & if/whether (Yes/No Questions)\\\hline \bt{als ob}& as if (mostly \bt{K2})\\\hline
        \bt{damit} & so that / in order to \\\hline \bt{\{so\}dass}& \{so\} that\\\hline \bt{wenn} & if\\\hline \bt{weil/da} & because\\\hline
        \bt{falls} & in case \\\hline \bt{ohne dass} & without\\\hline \bt{dadurch, dass} & due to the fact that\\\hline \bt{indem}& using\\\hline
        \bt{obwohl / obschon / obzwar / obgleich} & although \\\hline \bt{$\text{denn}^{\mathbb{N}}$} & because\\\hline
        \bt{$\text{deshalb}^{\mathbb{N}}$ / $\text{deswegen}^{\mathbb{N}}$ / $\text{daher}^{\mathbb{N}}$ / $\text{darum}^{\mathbb{N}}$} & hence\\\hline
        \bt{$\text{trotzdem}^{\mathbb{N}}$ / $\text{dennoch}^{\mathbb{N}}$}& nevertheless\\\hline
        \bt{wohingegen} & against which / whereas \\\hline \bt{bevor / ehe} & before\\\hline \bt{bis}& until\\\hline \bt{während} & while\\\hline
        \bt{solange} & as long as\\\hline \bt{sobald}& as soon as\\\hline \bt{sooft} (pronounced \bt{so/oft})& as often as / whenever \\\hline
        \bt{nachdem} & after\\\hline \bt{seit / seitdem} & since\\\hline \bt{Je\ldots desto/umso\ldots} & the\ldots the\ldots (comparators)\\\hline
    \end{tabular}
    \caption{Relationals: $\mathcal{R}_\boxdot$}
    \label{tab:destceT2}
\end{table}

The comparators \bt{Je\ldots desto/umso\ldots.}  only take comparative adjectives as arguments:
\begin{example}
    \begin{tabular}{m{7cm}|m{7cm}}
        \bt{Je [länger] ich auf dieser Erde lebe, desto [stärker] werde ich} & The longer I live on this earth, the stronger I will become
    \end{tabular}
\end{example}
\bt{Dadurch, dass} literally means `through the fact that', and the first part \bt{dadurch} doesn't have any consistent placement rules, at least not to my knowledge. It need only exist somewhere in the first argument. To check the nature (shifting/normal) of any unknown adverbial, simply use \textit{dwds.de}: a resource linked in \autoref{sec:dermrk}.

Next  are  conjunction contractions related to \bt{dass} and \bt{damit}. These are only allowed when the case $1$ nouns of both argument clauses  are the same.
\subsubsection{Dass - zu + Grundform}
Works for anything  that uses \bt{dass} as a component. Eliminate the case $1$ noun and conjunction \bt{dass} from the affected clause, and use $\vm{T}(:,4)$, with the CV being pushed to the end:
\begin{syntax}
    DC + \bt{zu} + CV (\bt{Grundform})
\end{syntax}
The special rules pertaining to $\mathcal{S}$ still apply when the CV is exiled. \bt{Zu} always appears before the very last verb in the chain. This contraction  represents an aspect, which is applied to the time  carried over from the main clause:
\begin{example}
    \begin{tabular}{m{6.5cm}|m{6.5cm}}
        \textit{Expanded}& \textit{Contracted}\\\hline
        Der Mann ist  klug genug, dass er seinen Weg   herausfindet  & Der Mann ist klug genug, seinen Weg  herauszufinden \\\hline
        Er konnte nicht glauben, dass er verloren hat & Er konnte nicht glauben, verloren zu haben\\\hline
        Er konnte nicht glauben, dass er so viel  verlieren kann & Er konnte nicht glauben, so viel  verlieren zu können\\\hline
        Der Mann ist glücklich, dass er seinen Weg  hat herausfinden können & Der Mann ist glücklich, seinen Weg   herausgefunden haben zu können
    \end{tabular}
\end{example}
The argument with `that' either occurs along/after  its other argument, or is already complete before it. Hence, ONLY cells  $\vm{T}(:,2)$ can be used in the argument clause with \bt{dass}, which is why only two contraction types using $\vm{T}(:,4)$ exist and suffice. All cells can otherwise be used in the other argument. 

\textit{An exception} occurs with the $\vm{T}(2,4)$ type contraction if any $V\in\mathcal{S}$ is involved. In this case, the main verb is no longer permeable and interacts with the exiled $R\in\mathcal{X}$, causing a $\vm{T}(2,2)$ chain to be formed. Note that $R$ must follow the main verb due to this interaction according to $\vm{T}(2,2)$. All other rules apply as normal.

Sometimes, if the main clause has a small DC,  the contracted argument is placed before it as part of the sentence:
\begin{example}
    \begin{tabular}{m{7cm}|m{7cm}}
        \bt{Ich werde Besseres [mit meinem Geld anzufangen] wissen} & I will know to start something better with my money
    \end{tabular}
\end{example}
\subsubsection{Damit - um\ldots zu + Grundform}
Replaces the adverbial \bt{damit} meaning `in order to/so that'. The rules are exactly the same as the above \bt{zu + Grundform}, except that now an extra \bt{um} stands at the beginning of the contracted clause:
\begin{example}
    \begin{tabular}{m{7.5cm}|m{7.5cm}}
        \textit{Expanded}& \textit{Contracted}\\\hline
        Er verwendete die Segen der Dämonen, damit er sich aus der Schwierigkeit aufheben konnte& Er verwendete die Segen der Dämonen, um sich aus der Schwierigkeit aufheben zu können\\\hline
        Er zählt die Schafe, damit er die Winterzeit aushält & Er zählt die Schafe, um die Winterzeit auszuhalten
    \end{tabular}
\end{example}
$\vm{T}(2,4)$ based constructions are  rare for this contraction type due to its meaning.
\subsubsection{Ohne zu}
This shouldn't technically be here but there is no other fitting place for it:
\begin{setdef}
    \begin{tabular}{m{7.5cm}|m{6.5cm}}
        \textit{Rule 1} & \textit{Rule 2}\\\hline
        Verbs of perception (\bt{sehen, hören\ldots}) & \bt{bleiben, sitzen, stehen, liegen}\\\hline
        \bt{helfen} & Verbs of motion (\bt{gehen, fahren, kommen\ldots})\\\hline
        - & \bt{lernen, lehren, üben} 
    \end{tabular}
\end{setdef}
Under such usage, the aforementioned verbs take verb arguments according to the rules: 
\begin{enumerate} 
    \item Temporarily treated as members of $\mathcal{S}$ 
    \item Treated similar to $\mathcal{S}$, with the only difference being the \bt{Perfekt} aspect.  \bt{Perfekt}ion  occurs normally, and there are no double infinitive rules as otherwise observed with  members of $\mathcal{S}$
\end{enumerate}
Verb arguments serve as meaning supplements. There is unfortunately no good explanation as to why  $\mathcal{S}$ allows the second category of verbs only partial citizenship.
\begin{example}
    \begin{tabular}{m{6.5cm}|m{7cm}}
        \bt{Wir wollen hier stehen bleiben} & We want \ut{here standing stay}\\\hline
        \bt{Sie ist heute früher spazieren gegangen} & She has \ut{today earlier} \ut{walking gone}\\\hline
        \bt{Er hat ihn rennen sehen} & He has \ut{him} \ut{running seen}
    \end{tabular}
\end{example}
Original members from $\mathcal{S}$ can also be used on such chains, as illustrated in the first example.

\subsection[Interrogation]{Interrogation (Fragesatz)}
$\biguplus$ We start with the \textit{Yes/No} (\bt{Ja/Nein}) questions:
\begin{syntax}
    Switch the  case $1$  and CV units
\end{syntax}
Practically, this often means switching the first two units, so that the CV unit appears  first. Since conjunction usage is invalid in questions, the first clause of the SPR always holds, thereby simplifying this formation:
\begin{example}
    \begin{tabular}{c|c}
        \bt{Bist du tot?} & Are you dead?\\\hline
        \bt{Wollen wir uns wieder treffen?} & \ut{Want we us again meet?}\\\hline
        \bt{Ist er zu früh angekommen?} & Has he too early \ut{arrived}?
    \end{tabular}
\end{example}
The interrogative words (IWs) used for \textit{elaborations} are as follows:
\begin{setdef}
    \begin{tabular}{c|c|c|c}
    \bt{Wann} & When & \bt{Wie} & How\\\hline \bt{Wo} & Where & \bt{Was}& What\\\hline
    \bt{Warum}& Why &\bt{Wieso} & Why\\\hline \bt{Wer*} & Who & \bt{Welcher*}& Which
    \end{tabular}
\end{setdef}
The words with the * change form, while all others are invariant. The form variation of \bt{welcher} is exactly as in \autoref{tab:denounT9}, whereas \bt{wer} follows \bt{der} ($\vm{Q}(:,1)$). 

\bt{Warum} focuses on the underlying reason behind an occurrence, whereas \bt{wieso} asks how  the occurrence came to pass:
\begin{example}
    \begin{tabular}{c|c|c}
        \textit{Question}& \bt{Warum} & \bt{Wieso}\\\hline
        Why did you fall?& The floor was wet & I slipped on the wet floor
    \end{tabular}
\end{example}
Using one of these IWs as the first unit in a sentence declares an elaboration:
\begin{example}
    \begin{tabular}{c|c}
        \bt{Warum tust du das?}& Why are  doing \ut{you} this?\\\hline
        \bt{Woher kommst du?} & From where \bt{come} you?\\\hline
        \bt{Wohin willst du?} & \ut{To where} \ut{want} you (go)?
    \end{tabular}
\end{example}
The situation is slightly different with \bt{wie}, which often deals with scales and properties, thus requiring adverbs. Interestingly, \bt{wie} and the adverb  are treated as a single unit for the SPR:
\begin{example}
    \begin{tabular}{m{6.5cm}|m{7cm}}
    \bt{Wie spät ist es?} & How \ut{late} is it?\\\hline
    \bt{Wie bist du hierhergekommen?} & How have you come here?\\\hline
    \bt{Wie schnell  fahren wir?} & How fast are \ut{driving} we?
    \end{tabular}
\end{example}
As a reminder, adverbs do not have  gradations, so \bt{Wie schneller} is not allowed.
\subsection[Describer subclause]{Describer subclause (Beschreibungssatz)}
$\biguplus$ \autoref{tab:destceT3} represents $\vm{Q}$ with all  hook word forms. The syntax of German D-subclauses is exactly the same as English: Commas separate the subclause, which is introduced using a hook word, from the main sentence. If a preposition is used, it is placed exactly before the hook word. The second SPR clause is always applicable:
\begin{example}
    \begin{tabular}{m{7cm}|m{7cm}}
        \bt{Der Mann, der ein König geworden war, wird\ldots}& The man, who had  become king, is becoming\ldots\\\hline
        \bt{Das Kind, dessen Haustier allmählich sterben wird, war\ldots} & The child, whose pet will be dying gradually, was\ldots\\\hline
        \bt{Die Frau, deren Kind einen Wutanfall gemacht hatte, ist angekommen} & The woman, whose child had thrown a tantrum, has arrived
    \end{tabular}
\end{example}
$\vm{Q}$ can also be used as an emphasizing replacement for  $\vm{P}[a,In/De,3_p]$. This feature is formally introduced in \autoref{sec:despco},  and requires clear prior context for use:
\begin{example}
    \begin{tabular}{m{7cm}|m{7cm}}
        \bt{Einführung in elektrische Systeme und deren Anwendungen\ldots} & Introduction to electrical systems and their (the electrical systems specifically) applications\\\hline
        \bt{Einführung in elektrische Systeme und ihre Anwendungen} & Introduction to electrical systems and their (could be something mentioned earlier, maybe mechanical systems) applications
    \end{tabular}
\end{example}
Pronoun based subclauses require iteration of the pronoun, since $\vm{Q}$ strictly assumes $3_p$:
\begin{example}
    \begin{tabular}{m{7cm}|m{7cm}}
        \bt{Ich, der ich dieses Weltall erstellt habe, bin gelangweilt}& I, who \ut{this universe} \ut{created have}, am bored\\\hline
        \bt{Du, die du mit dem Kind gekämpft hast, bist trottelig} & You, who have \ut{with the child} \ut{fought}, are moronic
    \end{tabular}
\end{example}
The gender of $\vm{Q}$ is taken from the real gender of the person, similar to adjective based nominalizations.
\begin{table}[bth]
    \centering
    \begin{tabular}{|c|c|c|c|c|}
        \hline$\diamond$ & \bt{Männlich} & \bt{Weiblich} & \bt{Sächlich}& \bt{Mehrzahl}\\\hline
        \bt{Nom} & der & die & das & die\\\hline
        \bt{Akk} & den & die & das & die\\\hline
        \bt{Dat} & dem & der & dem & denen \\\hline
        \bt{Gen} & dessen & deren & dessen & deren\\\hline
    \end{tabular}
    \caption{Hook words: $\vm{Q}$}
    \label{tab:destceT3}
\end{table}

For indefinite descriptions, the IWs: \bt{was, welcher} are used with their corresponding form variations. \bt{Welcher} is used only for specific references in technical German, similar to the replacement possessives as described above. \bt{Was} is used as usual:
\begin{example}
    \begin{tabular}{m{6.5cm}|m{7cm}}
        \bt{Das Haus, welches sie gekauft haben, ist alt} & The house, which you \ut{bought have}, is old \\\hline
        \bt{Nichts, was wir machen, ist umsonst}& Nothing \ut{what} we do is \ut{for nothing}\\\hline
        \bt{Das Beste, was ich je gesehen habe} & The best \ut{what} I ever seen \ut{have}
    \end{tabular}
\end{example}
Contracted adjectives and sentence units also require \bt{was}, if further described. For the CV, \bt{was} is always treated as a third person singular, if used in case $1$. Under certain circumstances, \bt{warum, wo} can also be used as hook words, but these add  sentence units  rather than actual describer subclauses:
\begin{example}
    \begin{tabular}{m{7cm}|m{7cm}}
        \bt{Das Haus, wo ich wohne\ldots} & The house, where I live\ldots\\\hline
        \bt{Das Haus, in dem ich wohne\ldots}& The house, in which I live\\\hline
        \bt{Der Grund, warum ich nicht kommen kann\ldots} & The reason why I cannot come\ldots\\\hline
        \bt{Der Grund, wegen dessen ich nicht kommen kann\ldots} & The reason, because of which I cannot come
    \end{tabular}
\end{example}
\subsubsection{MdG Verkürzung}
The \bt{Mittelwort der Gegenwart (MdG)}, alternatively known as \bt{Partizip1} is an analogue to the English P1, which is formed by adding a \bt{-d} to the \bt{Grundform}. The sole exception is \bt{sein/seiend}, which doesn't end up mattering because it has no use case. None.
\begin{example}
    \begin{tabular}{m{7cm}|m{6cm}}
    Er fährt mit dem Bus, der hell und heiß brennt, zur Arbeit & Er fährt mit dem (hell und heiß brennenden) Bus zur Arbeit\\\hline
    Sie aß sein Bein, das schmerzhaft ausblutete, voller Freude & Sie aß sein (schmerzhaft ausblutendes) Bein voller Freude
    \end{tabular}
\end{example}
\bt{Voller Freude} is an old German relic representing case $6$, which means `full of joy'. Covered later in \autoref{sec:despco}. 

\bt{MdG} contractions can also be carried out for the \bt{sein + zu + Grundform} passive replacement. The CV is \bt{sein}, and \bt{zu + Grundform} immediately precedes it. For the contraction, \bt{sein} is eliminated, while the \bt{Grundform} becomes the \bt{MdG}:
\begin{example}
    \begin{tabular}{m{7cm}|m{7cm}}
        \textgerman{Wir werden morgen die frühen Prüfungen, die unter \enquote*{Altklausuren} zu finden sind, behandeln} & \textgerman{Wir werden morgen die frühen (unter \enquote*{Altklausuren} zu findenden) Prüfungen behandeln}
    \end{tabular}
\end{example}
There is no \bt{zu + MdV} counterpart to this contraction type.
\subsubsection{MdV Verkürzung}
The \bt{Mittelwort der Vergangenheit (MdV)}, alternatively known as \bt{Partizip2} is an analogue to the English P2, whose formation has been dealt with in \autoref{sec:deverb}.

This is the only point where the diligent reader may get confused, because traditional German grammar and I treat an \bt{MdV} expansion differently. Although I consider it as the \bt{Zustandsleidform} (state passive), grammar textbooks often use the \bt{Vorgangsleidform} (process passive) instead. Clearly, a justification is in order.

An embedded subclause acts identical to an adjective, in the sense that it describes a property of the associated noun. Much more important than the process of how this property came to be, is the fact that it exists. English agrees on this respect, and the consequent expansion is  logical and easier to grasp. The \bt{Vorgangsleidform} is quite counterintuitive and does not fit in many cases, since the `Ongoing' aspect references that the action is currently in progress at the time of reference. On the other hand, the aspects do not mean exactly what one may expect for the \bt{Zustandsleidform}, which will be clarified in the section on passives, and make the subclause interpretation intuitive.

Although one may argue that the \bt{Perfekt} aspect with process passives could be used, standard grammar does not allow this, and the rules surrounding subclause contractions are strict to begin with. Hence, any expansion of \bt{MdV} subclauses  will explicitly use the \bt{Zustandsleidform} in this solution, which does not follow standard texts:
\begin{example}
    \begin{tabular}{m{7cm}|m{7cm}}
        Die Frau, die trotz fehlender Beteiligung in der Masche getötet ist, ist schön & Die (trotz fehlender Beteiligung in der Masche getötete) Frau ist schön\\\hline
        Den Aufbau der Leistungselektronik, die in den verschiedenen Fahrzeugsteuergeräten eingesetzt waren, \ldots & Den Aufbau der (in den verschiedenen Fahrzeugsteuergeräten eingesetzten) Leistungselektronik\ldots
    \end{tabular}
\end{example}
Although the given examples do not reflect this, embedding is independent of other noun unit features: \bt{schnell fahrendes Auto / geschicktes Bild} without any articles is valid.

\ut{\bt{Warning}}: \bt{MdV} contractions cannot be  performed for any verb $V\in\mathcal{S}$.
\subsubsection{Miscellaneous}
\begin{enumerate}
\item The modal verbs \bt{Müssen, Können} can  be replaced using \bt{zu + Grundform}, which falls under the category  `Modal replacements'. Among them, \bt{müssen} may only be replaced in the active voice:
\begin{setdef}
    \begin{tabular}{m{4cm} c}
        \textit{Modal Verb} & \textit{Replacement}\\\hline
        Müssen  & haben + zu + Grundform \\\hline
        Können & Haupttätigkeitswort + [zu + Grundform][\bt{D-Subclause}]
    \end{tabular}
\end{setdef}
The first verb is the CV. Since  \bt{können}  uses $\vm{T}(1,4)$ for contraction, only $\vm{T}(1,:)$ is allowed in the subclause:
\begin{example}
    \begin{tabular}{m{7cm}|m{7cm}}
        \textit{Contracted} & \textit{Expanded}\\\hline
        Ich habe Brot zu schneiden & Ich muss Brot schneiden \\\hline
        Sie fand nichts zu sagen & Sie fand nichts, was sie sagen konnte\\\hline
        Ich habe ihn zu töten gehabt & Ich habe ihn töten müssen \\\hline
        Es gibt nur wenig zu machen & Es gibt nur wenig, was gemacht werden kann
    \end{tabular}
\end{example}
\item Both \bt{MdG, MdV} describer subclause contractions can be written in a different form, if a noun unit embedding is not desired. Simply eliminate the hook word and CV (assuming the DC is not a null set), and let the remaining phrase stay as is:
\begin{example}
    \begin{tabular}{m{7cm}|m{6.8cm}}
        \textit{Expanded} & \textit{Contracted}\\\hline
        Die Braut, die vor Glück strahlte, sah in die Kamera & Die Braut, vor Glück strahlend, sah in die Kamera\\\hline
        Ihr Vater, der von Wehmut ergriffen war, wandte sich ab & Ihr Vater, von Wehmut ergriffen, wandte sich ab
    \end{tabular}
\end{example}
This is especially useful if the noun is the name of a person. Sometimes, if the main sentence is small, the shortened subsentence can also be thrown to the very end, though this might cause confusion and should be avoided:
\begin{example}
    Ihr Vater wandte sich ab, von Wehmut ergriffen
\end{example}
\item Pronoun noun duplication is possible in general  for $\vm{P}[a,In,:](1,\mathbb{M},0\leftrightarrow\mathbb{H})$:
\begin{example}
    \begin{tabular}{c c|c c}
        \bt{Wir Reisenden\ldots} & We the travelling\ldots & \bt{Ihr Suchenden\ldots} & You the seeking\ldots\\\hline
        \bt{Wir Gefangenen\ldots} & We the prisoners\ldots & \bt{Sie Verletzten\ldots} & They the hurt\ldots\\\hline
        \bt{Sie Ungeheuer\ldots} & You monsters\ldots & \bt{Ihr Fahrer\ldots} & You drivers\ldots
    \end{tabular}
\end{example}
This may also  be done for some $\vm{P}[a,In,:](1,1,0\leftrightarrow\mathbb{H})$ under specific conditions: \bt{du Süße}
\item \bt{Wenn}-arguments using cell $\vm{T}(1,2)$ with \bt{man, es} as the case $1, 2$ nouns respectively may also be contracted. To do this, \bt{wenn, es} and \bt{man} are eliminated, and the CV is converted into its \bt{MdV} form. Sometimes the unit ordering can also be flipped:
\begin{example}
    \begin{tabular}{m{7cm}|m{7cm}}
        \textit{Contracted} & \textit{Expanded}\\\hline
        genau genommen & Wenn man es genau nimmt\ldots\\\hline
        verglichen mit [Zitronen] & Wenn man es mit [Zitronen] vergleicht\ldots\\\hline
        abgesehen davon & Wenn man es davon absieht\ldots\\\hline
    \end{tabular}\\
    anders ausgedrückt/gesagt, angenommen, gesetzt [in] den Fall, kurzgesagt, so gesehen/betrachtet, vorausgesetzt\ldots
\end{example}
\item Noun duplication  follows the typical \{\textit{CoM} content \textit{CoM}\} pattern:
\begin{example}
    Die Quelle [des Rheins], [des längsten Flusses innerhalb Deutschlands], liegt in der Schweiz
\end{example}
\end{enumerate}

\subsection[Order/Request]{Order/Request (Befehl/Aufforderung)}
\begin{setdef}
    \begin{tabular}{c|m{7cm}}
        Order cell & $\vm{T}(1,2)$ \\\hline
        Subsubsection naming & $\vm{S}[\vm{T}(1,2)](2_p,m,h)$, where ($m,h$) are variables\\\hline
        Grade keyword & \bt{Bitte}\\\hline
        Hypothesis based cmd & \bt{Möglichkeitsform (K2)} of \bt{können / haben}
    \end{tabular}
\end{setdef}
Compounding grade keywords and hypotheticals cannot cross into the next integral range, unless the grade is physically increased. \bt{Haben} hypotheticals are treated separately at the end.
\subsubsection{Cmd:(1,0)}
\begin{enumerate}
    \item The CV is kept as the first unit
    \item Eliminate any \bt{Umlaut} that appears due to irregular conjugation, and the \bt{-st} end. Sometimes, an additional \bt{-e} may appear. (Do not eliminate natural \bt{Umlaut})
    \item Eliminate \bt{du}, and add an \textit{ExC} (!), \textit{InQ} (?) or \textit{EoS} (.) at the end
    \item Add any grade keywords if desired
\end{enumerate}
Hypothesis commands follow interrogative structures, so only the last step  is applicable for them.
\begin{example}
    \begin{tabular}{m{6cm}|m{6cm}|c}
        \textit{Command} & \textit{Original/Translation} &\textit{Grade}\\\hline
        Könntest du bitte gehen? & \bt{Could you please go?} & $\approx0.8$\\\hline
        Töte die Ungerechtigkeit der Welt, bitte& Du tötest die Ungerechtigkeit der Welt & $0.5$\\\hline
        Zieh dich an! & Du ziehst dich an  & $-0.5$\\\hline
        Iss nicht normal & Du isst nicht normal & $0$
    \end{tabular}
\end{example}
Exceptions: \bt{sein/sei}, \bt{haben/hab}. 
\subsubsection{Cmd:(M,0)}
\begin{enumerate}
    \item The CV is kept as the first unit
    \item Eliminate \bt{ihr}, and add an \textit{ExC} (!), \textit{InQ} (?) or \textit{EoS} (.) at the end
    \item Add any grade keywords if desired
\end{enumerate}
\begin{example}
    \begin{tabular}{m{6cm}|m{6cm}|c}
        \textit{Command} & \textit{Original/Translation}& \textit{Grade}\\\hline
        Geht heim! & \bt{Go home!}& $-0.5$\\\hline
        Entschuldigt euch für dieses Verhalten & \bt{Excuse yourselves for this behaviour} & $0$\\\hline
        Könntet ihr leise sein? & \bt{Could you  quiet \ut{be}} &$0.5$
    \end{tabular}
\end{example}
\subsubsection{Cmd:(\{0,M\},H)}
\begin{enumerate}
    \item The CV is kept as the first unit and \bt{Sie} as the second
    \item Add an  \textit{InQ} (?) or \textit{EoS} (.) at the end
    \item Add any grade keywords if desired
\end{enumerate}
The highest politeness level does not accept exclamations, for obvious reasons.
\begin{example}
    \begin{tabular}{m{6cm}|m{6cm}|c}
        \textit{Command} & \textit{Original/Translation}& \textit{Grade}\\\hline
        Kommen Sie bitte herein & \bt{\ut{Come} \ut{You} \ut{please} inside} & $\mathbb{H}+0.5$\\\hline
        Entspannen Sie sich & \bt{Relax \ut{You} \ut{yourself}}& $\mathbb{H}$\\\hline
        Könnten Sie bitte einen Platz nehmen? & \bt{Could you please a seat \ut{take}?}& $\approx\mathbb{H}+0.8$
    \end{tabular}
\end{example}
Exception: \bt{sein/seien}

Rewriting \bt{Sie} based commands that do not employ hypotheticals can create signboard instructions. Simply use the second SPR clause and eliminate \bt{Sie}:
\begin{example}
    \begin{tabular}{m{7cm}|m{7cm}}
        \textit{Command} & \textit{Signboard}\\\hline
        Benutzen Sie (den) Aufzug im Brandfall nicht & Aufzug im Brandfall nicht benutzen\\\hline
        Schütteln Sie (den Saft) gut vor dem Öffnen & Vor dem Öffnen gut schütteln
    \end{tabular}
\end{example}
Often, the case $2$ unit is omitted or used without any additional features.
\subsubsection{Miscellaneous}
The \bt{haben K2} form is essentially a normal active voice \bt{K2} sentence in $1_p$. Without going into the technicalities of \bt{K2}, here are some examples. Note that this is still $\vm{T}(1,2)$:
\begin{example}
    \begin{tabular}{m{7cm}|m{7cm}}
        \bt{Ich hätte gern eine Banana} & I would like to have a banana\\\hline
        \bt{Wir hätten gern eine Fahrkarte} & We would like to have a ticket
    \end{tabular}
\end{example}
Mostly useful during shopping. The last remaining command form is the $1_p$ group commands, which follow the English structure using `let':
\begin{example}
    \begin{tabular}{c|c}
        \bt{Lass uns gehen!} & Let us go!\\\hline
        \bt{Los geht es!} & \ut{Off goes it}!
    \end{tabular}
\end{example}
The \bt{Cmd:(1,0)} form is applied on \bt{Lassen} for such group orders. \bt{Los geht's} is also commonly used as a German specific phrase.

\subsection[Hypothesis]{Hypothesis (K2/Möglichkeitsform)}
$\biguplus$ First, $\vm{T}_K$ must be defined. Only the aspect `Ongoing' exists, with the following distribution:
\begin{equation*}
    \vm{T}_K(1,1)=\vm{T}(2,2), \quad \vm{T}_K(1,2)=\vm{T}(1,2), \quad \vm{T}_K(1,3)=\vm{T}(1,3)
\end{equation*}
In order to apply the mode \bt{K2} on any CV and get its \bt{K2} root, the following steps are necessary:
\begin{enumerate}
    \item Take the \bt{Präteritum} root: $\vm{T}(1,1)$ of the CV
    \item If an \bt{Umlaut} can be added to an existent \bt{a/o/u}, add it
\end{enumerate}
Hence, $\vm{T}_K(1,1)\neq\vm{T}(1,1)$  to avoid overlap with $\vm{T}_K(1,2)$. That being said, the overlap is not completely absent, since a lot of verbs do not have any \bt{Umlaut}able letters, which makes the \bt{K2} and \bt{Präteritum} forms  coincide. Also, some verbs do not allow an \bt{Umlaut} addition despite having \bt{Umlaut}able letters.

To avoid this confusion, the $\vm{T}_K(1,3)$ form is often used as a substitute for $\vm{T}_K(1,2)$, as future assumptions are rare anyway. CV form variation follows the pattern in $\vm{T}(1,1)$, but with the \bt{K2} root.

The conjunction \bt{wenn} is used often when speaking of conditions, and is the only one which allows \bt{K2} usage within its arguments. The only exceptions appear for the sentence mode \bt{mittelbare Wirkart}, where \bt{dass, ob} among others may also take \bt{K2} arguments.

\bt{Wenn} is not strictly necessary for \bt{K2} usage - double usage  simply increases the uncertainty of the considered statements:
\begin{example}
    \begin{tabular}{m{7cm}|m{6.8cm}}
        \bt{Ich würde mich freuen, wenn er gehen würde}& I would \ut{myself} \ut{delight}, if he \ut{go} \ut{would}\\\hline
        \bt{Wenn ich [bloß/doch] ein Millionär wäre!} & If \ut{I} [only] a millionare \ut{were}!\\\hline
        \bt{Wären sie früher angekommen, wären wir früher gegangen} & Had \ut{they} earlier \ut{arrived}, would \ut{we} earlier \ut{have gone}\\\hline
        \bt{Wenn Sie  uns früher Bescheid gegeben hätten, hätten wir früher kommen können} & If you had earlier \ut{notice} \ut{given}, would \ut{we} \ut{earlier} have been able to come\\\hline
        \bt{Nähmest du einen Schritt weiter, sähest du seinen Tod} & If you take  a step further, you will see his death\\\hline
        \bt{Ginge die Sonne auf, käme die Vögel} & If the sun \ut{goes up}, the birds would come
    \end{tabular}
\end{example}
There are two special constructions associated with \bt{K2}:
\subsubsection{zu + als dass}
Too (something) to (something). Don't take  \bt{als dass}  word by word, it just means `to' collectively:
\begin{example}
    \begin{tabular}{m{7cm}|m{7cm}}
        \bt{Dieses Blut ist zu lecker, als dass ich es verpassen würde}  & This blood is too delicious to miss \\\hline
        \bt{Die Sanierung war zu teuer, als dass sie sich gerechnet hätte} & The renovation was too expensive to be cost effective
    \end{tabular}
\end{example}
The \bt{zu} in the first clause is necessary: it represents an exaggeration that due to the first condition, the second is highly improbable. The \bt{K2} time is matched according to the first clause.
\subsubsection{als ob / als wenn / als}
This is a direct analogue to `He behaves, as if\ldots' Sometimes, you can find \bt{K1} forms used for this purpose in written language:
\begin{example}
    \begin{tabular}{m{7cm}|m{6.5cm}}
        \bt{[Er tut so], als ob er Gott wäre} & [He pretends] as if he were god\\\hline
        \bt{[Du siehst aus], als hättest du tagelang durchgearbeitet} & [You look] like you have worked throughout the entire day
    \end{tabular}
\end{example}
Here, the \bt{K2} time is independent of the introductory clause.

\subsection[Voice]{Voice (Wirkart)}
$\biguplus$ The active voice will be termed \bt{Tätigform} while the passive as the \bt{Leidform}. The subject suppliers are:
\begin{setdef}
$\circledS = \begin{bmatrix}
    \texttt{von} & \texttt{durch}
\end{bmatrix}$
\end{setdef}
The process and state passives work exactly as described in \autoref{sec:enstce}. The process passive places emphasis on the fact that the action occurred, whereas the state passive is concerned only with the consequences of the action.
\subsubsection{Vorgangsleidform}
Also called the \bt{werden} ($V\in\mathcal{X}$) passive. To obtain its base cell:
\begin{enumerate}
    \item Send the \bt{MdV} form of the CV to the DC
    \item The new CV is \bt{werden}
    \item \bt{\ut{Exception}}: Iff  the \bt{Perfekt} form is necessary,  \bt{werden}  will become \bt{worden} in the DC instead of \bt{geworden}
\end{enumerate}
This creates $\vm{T}(1,2)$ for the process passive, which enables sentences like:
\begin{example}
    \begin{tabular}{m{7cm}|m{6cm}|c}
        \textit{Statement} & \textit{Translation}& \textit{Cell}\\\hline
        \bt{Das muss gemacht werden} & This must be done & $\vm{T}(1,2)$\\\hline
        \bt{Sein Feind ist  getötet worden} & His enemy has  been killed & $\vm{T}(2,2)$\\\hline
        \bt{Das Schwert wird gerostet worden sein müssen, als wir zurückkehren} & The sword will have to have been rusted, when  we return & $\vm{T}(2,3)$
    \end{tabular}
\end{example}
Note that since \bt{worden} is not a modal or main verb, it blocks any exiled $K\in\mathcal{X}$ from advancing leftwards beyond it, according to the second special rule for any $V\in\mathcal{S}$. 

By virtue of its meaning, the process passives for \bt{wollen} are formed by substituting it with \bt{sollen}:
\begin{example}
    \begin{tabular}{m{7cm}|m{7cm}}
        Man will einen neuen Kindergarten aufbauen & Ein neuer Kindergarten soll aufgebaut werden
    \end{tabular}
\end{example}
No other $V\in\mathcal{S}$ is substituted like this.

Although the process passive could technically be formed for reflexive usage, this is not grammatically allowed, and any occurrences are strictly colloquial. A frequently observed  use concerning the \bt{Vorgangsleidform} is the implication of $\vm{T}(1,3)$ despite using $\vm{T}(1,2)$:
\begin{example}
    \begin{tabular}{m{8cm}|m{6cm}}
        $\vm{T}(1,3)$ \textit{Statement} & \textit{Translation}\\\hline
        \bt{Widerrechtlich abgestellte Fahrzeuge werden abgeschleppt} & Illegally parked vehicles will be  towed\\\hline
        \bt{Die Gäste werden später begrüßt} & The guests will be greeted later
    \end{tabular}
\end{example}
Such examples depict situations where the action is fixed, rulebound or supposed to occur immediately in the near future. These conditions mimic the `simple' aspect  that English uses, which is why the best translation uses the English cell - $\vm{T}(1,3)$. However, since \bt{Deutsch} does not possess this aspect, the above method is a  workaround to create $\vm{T}(\texttt{simple},3)$. There exist no such workarounds for any other cells or sentence modes.

To reduce the DC length  when a $V\in\mathcal{S}$ is involved, there exist alternative constructions for  the \bt{Vorgangsleidform}.  However, not all $V\in\mathcal{S}$ can be replaced by any given method.

Since there are many such possible reformulations, only the 4 most important  have been listed below. This exclusion is inconsequential, since all of the omissions are based either on rule `1' or `2', and can replace either \bt{müssen} or \bt{können}:
\begin{setdef}
    \begin{tabular}{m{0.5cm}|m{6cm}|m{5cm}}
        $\clubsuit$ & \textit{Replacement}&  $V\in\mathcal{S}$\\\hline
        1 & Sein + zu + Grundform  & Müssen, Können, Sollen, Nicht dürfen\\\hline
        2 & Lassen sich + Grundform  & Können\\\hline
        3 & Sein + Eigenschaftswort(-bar/-lich) & Können\\\hline
        4 & Man  & Keine
    \end{tabular}
\end{setdef}
The fourth pattern is neither passive nor a replacement, but the performer is still unknown, and thus retains a passive-like meaning despite being in active voice. Examples for $\vm{T}(1,2)$:
\begin{example}
    \begin{tabular}{m{6.5cm}|m{7cm}}
        \bt{Replacement} & \bt{Original}\\\hline
        Der Bericht ist zu schreiben & Der Bericht [kann / muss / soll] geschrieben werden\\\hline
        Der Bericht ist nicht zu schreiben & Der Bericht darf nicht geschrieben werden\\\hline
        Die Gleichung lässt sich schöner schreiben & Die Gleichung kann schöner geschrieben werden\\\hline
        Die Schrift ist gut lesbar & Die Schrift kann gut gelesen werden\\\hline
        Man baut ein Haus & Ein Haus wird gebaut 
    \end{tabular}
\end{example}
The CV and $\vm{T}$-cell calculations in each case should be obvious.  \bt{Zu + Grundform} is  an inert component that remains at the end of the sentence (beginning of the DC) and doesn't participate in cell transformation.

Lastly, although it should be obvious, it must be mentioned that any passive replacement (except 4) cannot be subject to another passive voice change. So
\begin{example}
Der Bericht wird zu schreiben gewesen : $\vm{T}(1,2)$
\end{example}
is absolute nonsense. Since 4 is technically active voice, it can be passived, which just converts the sentence back into its original form.

\subsubsection{Zustandsleidform}
Also called the \bt{sein} ($V\in\mathcal{X}$) passive. To obtain its base cell:
\begin{enumerate}
    \item Send the \bt{MdV} form of the CV to the DC
    \item The new CV is \bt{sein}
\end{enumerate}
The aspects of the \bt{Zustandsleidform} have a slightly different interpretation than normal. Since the state passive focuses on the consequence of an action, $\vm{T}(1,:)$ means that the consequence of an action currently exists. $\vm{T}(2,:)$ implies that the consequence existed, but no longer exists. Since the state passive declares  an  absolute and guaranteed state, no $V\in\mathcal{S}$ can be used with it.
\begin{example}
    \begin{tabular}{m{7cm}|m{7cm}}
        \bt{Die Tür ist geschlossen} & The door is closed\\\hline
        \bt{Die Tür ist geschlossen gewesen} & The door has been being closed
    \end{tabular}
\end{example}
The second example means that the door being closed was a condition that existed, but has since stopped existing. Although both perspectives seem to deal with this similarly, English is focusing on the action, while \bt{Deutsch} is looking at its consequence. $\vm{T}(2,:)$ for the \bt{Zustandsleidform} is practically never used.
\subsection[Subjective usage]{Subjective usage (Wahrscheinlichkeitsgebrauch)}
$\biguplus$ A verb $V$ may be used subjectively  under specific conditions:
\begin{enumerate}
    \item $V\in$ defined probability markers
    \item Special $\vm{T}_S$-cell forms are employed
\end{enumerate}
Such a usage attaches the  probability associated with $V$ to the state/action described by the sentence. The possible markers along with their approximate probabilities are:
\begin{setdef}
    \begin{tabular}{c|c|c|c}
        \textit{V} & \textit{Probability (\%)} & \textit{V} & \textit{Probability (\%)}\\\hline
        Müssen & $99$&  Müssen (K2) & $90$\\\hline
        Dürfen (K2) & $70$&   Können & $60$\\\hline 
        Können (K2) & $\leq50$ & Mögen & $\leq50$
    \end{tabular}
\end{setdef}
According to the surity of the declaration, the appropriate marker should be used. There are only two real $\vm{T}_S$-cells:
    \begin{equation*}
\vm{T}_S(1,1)=\vm{T}(2,2), \quad \vm{T}_S(1,2)=\vm{T}(1,2), \quad \vm{T}_S^*(1,3)=\vm{T}_K(1,3) 
    \end{equation*}
No future forms exist in this pattern, so any future uncertainties are declared using $\vm{T}_K(1,3)$.  Applying the subjective sentence mode occurs as follows:
\begin{enumerate}
    \item Exile the CV to the DC
    \item The new CV is $V$, which always uses its $\vm{T}(1,2)/\vm{T}_K(1,2)$ forms (depending on if it requires \bt{K2})
\end{enumerate}
Note that probability markers $\notin\mathcal{S}$, and should be treated  separately. Hence, they can also be applied on the \bt{Zustandsleidform}, even though $\mathcal{S}$ usage is illegal with it:
\begin{example}
    \begin{tabular}{m{8cm}|m{5.5cm}}
        \bt{Er könnte haben kommen müssen} & He might have had to come\\\hline
        \bt{Das mag in diesem Fall stimmen} & That may be true in this case\\\hline
        \bt{Die Tür muss geschlossen sein} & The door must be closed\\\hline
        \bt{Die Tür muss aufgebrochen worden sein} & The door must have been broken open
    \end{tabular}
\end{example}

\subsection[Speech]{Speech (Redeweise)}
$\biguplus$ \bt{Deutsch} allows unit reordering, unlike English:
\begin{example}
    \begin{tabular}{m{7cm}|m{7cm}}
        \multicolumn{2}{c}{\textit{Direct speech}}\\\hline
        \textgerman{\bt{Der Lehrer sagte, \enquote{Ihr müsst die Aufgabe heute vervollständigen.}}} & The teacher said, \enquote{You must complete the task today.}\\\hline
        \textgerman{\bt{\enquote{Aus dem Weg!}, schrie er}} & He screamed, \enquote{Out of the way!}
    \end{tabular}
\end{example}
\subsubsection{Mittelbare Redeweise}
The IACs for indirect speech:
\begin{setdef}
    \begin{tabular}{c|c|c}
        \textit{Nature/Mode} & \textit{IAC} & \textit{Comments}\\\hline
        \bt{Aussage} & \bt{dass} & -\\\hline
        \bt{Fragesatz (Ja/Nein)} & \bt{ob} & Opener always first\\\hline
        \bt{Fragesatz (Erläuterung)} & - & The question remains a sentence unit\\\hline
        \bt{Befehl} & \bt{dass} & Often contracted 
    \end{tabular}
\end{setdef}
The explicit sentence mode `Speech' corresponds to \bt{Konjunktiv 1 (K1)}, for which  $\vm{T}\neq\vm{T}_R$.

The diligent reader would have realized by now that any CV  has only two types of form variations, corresponding to cells $\vm{T}(1,2)$ and $\vm{T}(1,1)$.  Hence, any sentence mode must have  defined CV forms for both these cells to have the `complete' range at its disposal. But for some reason, the sentence mode \bt{K1}  has no defined $\vm{T}_R(1,1)$ forms, which makes it impossible to create $\vm{T}_R(:,1)$.

To circumvent this, all of the statements made about the past relative to the opener are grouped into $\vm{T}_R(2,2)$ or the \bt{Perfekt} aspect. This does introduce ambiguity about   past statements, but fatal errors are rather rare. 

$\vm{T}_R(1,2)$ is based on $\vm{T}(1,2)$ with the conditions:
\begin{enumerate}
    \item The \bt{K1} root is the VR
    \item $(3_p,1)$ follows $(1_p,1)$
    \item An additional \bt{e} is \ut{always} added at the end of the \bt{K1} root in $(2_p,:,0)$
    \item If the \bt{Grundform} ends in \bt{-eln}, then the \bt{e} is eliminated in the \bt{K1} root
    \item \bt{Grundform} in \autoref{tab:deverbT2} implies `VR + \bt{en}' irrespective of whether the infinitive end was \bt{-n/-en}
\end{enumerate}
The \bt{K1} root knows no  irregularity. So \bt{sein} becomes \bt{sei/seien}, making for an incredibly simple form variation\ldots until you realize that half of these forms are exactly the same as  $\vm{T}(1,2)$. This is especially a problem for regular verbs, since it defeats the purpose of \bt{K1}, leaving the interpretation up to context. To compensate for this, \bt{K2} is used instead but with a \bt{K1} meaning, which is a strange oversight in an otherwise structured perspective. Since this \bt{K2} has nothing to do with hypotheticals, it will be renamed to \bt{K3} instead.

\bt{K3} forms are exactly the same as \bt{K2}, which means that only the three cells of $\vm{T}_K$ exist for use, further compressing the expressive range. Despite the drop in accuracy, one can still get by:
\begin{example}
    \begin{tabular}{m{7cm}|m{7cm}}
        \bt{Der Lehrer befahl uns,  unsere Aufgabe zu vervollständigen} & The teacher ordered us  to complete our task\\\hline
        \bt{Er forderte, dass wir aus dem Weg gingen} & He demanded, that we  get out of the way\\\hline
        \bt{Er erklärt, (dass) er sei müde} & He is declaring that he is tired\\\hline
        \bt{Er fragte uns, ob wir frei haben} & He asked us if we have free time\\\hline
        \bt{Er fragte uns, wann wir frei haben} & He asked us when we have free time
    \end{tabular}
\end{example}

Indirect arguments in the replacement  forms for \bt{Aussage} do NOT use the optional mode \bt{K1}. Prepositions  placed before the speaker unit in the opener will be marked $\mathfrak{A}$, while those placed after it $\mathfrak{E}$:
\begin{setdef}
    \begin{tabular}{c|c|c|c}
        \textit{Preposition} & \textit{Placement} & \textit{Meaning} & \textit{Surity}\\\hline
        \bt{laut} &  $\mathfrak{A}$ & according to & Very High\\\hline
        \bt{gemäß} & $\mathfrak{E}, \mathfrak{A}$ & in accordance with & High\\\hline
        \bt{nach} & $\mathfrak{A}$ & according to & Medium\\\hline
        \bt{zufolge} &  $\mathfrak{E}$ & in consequence to & Full
    \end{tabular}
\end{setdef}
The English counterparts have been provided solely to give an idea for why all of these require the speaker unit to use the \bt{Dativ} case. One may also indicate the perceived certainty of  the reported statements using the appropriate preposition. 

If \bt{wie}[as]($\mathfrak{A}$) is employed, the speaker unit must  stay in case $1$:
\begin{example}
\begin{tabular}{m{6cm}|m{7cm}}
\textit{Contracted} & \textit{Original}\\\hline
[Gemäß einer Forschung] gehen\ldots & Die Forschung sagt, dass\ldots gehen\\\hline
[Wie verschiedene Zeugen bestätigten], zeigte die Ampel rot & [Verschiedene Zeugen bestätigten], dass die Ampel rot gezeigt habe\\\hline
[Den Angaben zufolge], ist er tot & [Die Angaben ausdrücken], dass er tot sei
\end{tabular}
\end{example}
\subsubsection{Miscellaneous}
The verbs \bt{sollen, wollen} can be used to recount the words of another without confirming their accuracy:
\begin{example}
    \begin{tabular}{m{6.5cm}|m{7cm}}
    \bt{Er soll ein erfolgreicher Anwalt sein} & [According to rumours/I have heard that] he is a successful lawyer\\\hline
    \bt{Er will ein erfolgreicher Anwalt sein} & [He claims/Supposedly] he is a successful lawyer
    \end{tabular}
\end{example}
The [text in square brackets] gives the  meaning encoded by \bt{wollen/sollen} in the respective cases. The  problem:  there are absolutely no VC differences with any regular uses that mean `want/should'. Instead, contextual clues do the heavy lifting, which is once again a strange oversight for an otherwise structured perspective.

Using  \bt{sollen, wollen} in this manner makes it possible  to  mix them with the \bt{Zustandsleidform}.

$\biguplus$ \bt{K1} is also used for the `Blessing' mode. Due to limited usage and the reuse of \bt{K1} forms, it does not deserve a separate subsection in \bt{Deutsch}:
\begin{example}
    \begin{tabular}{m{5cm}|m{5cm}|m{4cm}}
    \bt{Möge Gott euch segnen} & - & May god \ut{you} bless \\\hline
    \bt{Es lebe der König!} & \bt{Möge der König leben!} & May the king live!\\\hline
    \bt{Er fahre zur Hölle} & \bt{Möge er zur Hölle fahren} & May he go to hell\\\hline
    \bt{Gott helfe uns!} & \bt{Möge Gott uns helfen!} & May god help us!
    \end{tabular}
\end{example}
All the above sentences can  be universally written using \bt{mögen} based orders, but sometimes they are rephrased to sound natural/impactful.  \bt{Mögen} uses its subjective meaning in such usage.

The `Blessing' mode may also be used for generating indirect speech for a request:
\begin{example}
    \begin{tabular}{m{7cm}|m{8cm}}
        \bt{Der Kellner sagt, (dass) sie möge bitte Platz nehmen} & The waiter is requesting  her to take a seat\\\hline
        \bt{Die Abgeordneten baten, (dass) \{man\} möge ihrem Antrag zustimmen} & The representatives \ut{were requesting} that \{people\} consent to their application
    \end{tabular}
\end{example}

In technical \bt{Deutsch},  multiple considerations are often declared about something without hypothesizing: for example as the starting point for a proof, or derivation of a result, where it is more about logical reasoning and argumentation rather than unconfirmed statements. Such statements along `assumptions, argumentation, or instructions without a commanding tone'  use  \bt{K1}:
\begin{example}
    \begin{tabular}{m{5.5cm}|m{8cm}}
        \bt{Gegeben sei das skizzierte Rahmentragwerk} & \ut{Given is the sketched truss}\\\hline
        \bt{Betrachtet werde das Energiebordnetz eines Brennstoffzellenfahrzeugs} & The (on-board electrical system) of a (fuel cell automobile) is \ut{being} considered\\\hline
        \multicolumn{2}{m{13.5cm}}{\bt{Der Polizist rief dem Radfahrer zu, (dass) er  sofort anhalten müsse. [Er dürfe nicht auf die Autobahn fahren]} - [Instruction]}
    \end{tabular}
\end{example}
Since an assumption is to be declared, and the actor is not important, such structures are restricted to $\vm{T}_R(1,2)$ with the passive voice (unless dealing with instructions), and require the DC to stay as the first unit of the sentence. Such a technical restructuring is unique to \bt{Deutsch}.

\subsection[Negativity]{Negativity (Verneinung)}
$\vm{P}[b,In]$ based:
\begin{example}
    niemand, nichts, keiner, wenig, nirgendwo\ldots
\end{example}
The matrix $\vm{N}$ is given by:
\begin{setdef}
        $\vm{N} = \begin{bmatrix} \texttt{kein} & \texttt{nicht} \end{bmatrix}$
\end{setdef}
Noun unit negations are performed using \bt{kein}, which has been introduced as the negative article in \autoref{sec:denoun}:
\begin{example}
    \begin{tabular}{m{7cm}|m{7cm}}
    \bt{Das ist kein Spiel} & This is no game\\\hline
    \bt{Hier ist kein Geld zu verdienen} & There is no money to be earned here
    \end{tabular}
\end{example}
VC negations use  \bt{nicht}, which in general stands as  the first word in the unit it negates. If the meaning of the entire sentence is to be inverted, it stays at the very end, but before the DC:
\begin{example}
    \begin{tabular}{m{7cm}|m{7cm}}
        \bt{Sie haben mir das nicht erzählt} & You \ut{have} not told me this\\\hline
        \bt{So kannst du nicht überleben} & You can\ut{ }not survive like this\\\hline
        \multicolumn{2}{c}{\bt{Mengenangabe}}\\\hline
        \bt{Nicht jede Fruchtsorte bringt die gleichen Vorteile} & Not every  \ut{fruit type} brings the same advantages\\\hline
        \bt{Nicht alle verhalten sich vernünftig} & Not \ut{all} behave rationally
    \end{tabular}
\end{example}
\subsection[Timekeeping]{Timekeeping (Zeitmessung)}
$\mathfrak{R}24$:
\begin{syntax}
    (Hour) \bt{Uhr} (Minutes)
\end{syntax}
\begin{example}
    \begin{tabular}{c|c}
        acht Uhr fünfundfünfzig & 08:55\\\hline siebzehn Uhr zweiundvierzig & 17:42
    \end{tabular}
\end{example}
Markers for $\mathfrak{R}12$:
\begin{setdef}
    \begin{tabular}{c|c|c|c}
        \bt{halb} & half to (not past) & \bt{viertel} & quarter\\\hline
        \bt{vor} & before & \bt{nach} & after\\\hline
        \multicolumn{4}{c}{Uncertainty markers}\\\hline
        \bt{kurz} & shortly ($\approx5$min) & \bt{gleich / fast} & around / almost\\\hline
        \multicolumn{4}{c}{Half markers (HM)}\\\hline
        \bt{morgens} & AM & \bt{abends} & PM
    \end{tabular}
\end{setdef}
Only \bt{kurz} can be combined with \bt{vor, nach}. The reference time is  the future by default, with  time preposition  \bt{um}:
\begin{example}
    \begin{tabular}{m{6cm}|m{6cm}|c}
        \bt{Es ist kurz [vor] halb vier} & It \ut{is} shortly half past three & $\approxeq3:30-\epsilon$\\\hline
        \bt{viertel nach eins morgens} & quarter past one in the morning & $1:15$ AM\\\hline
        \bt{Es ist viertel vor zwei} & It is quarter to two & $1:45$\\\hline
        \bt{(kurz nach) zehn vor acht} & (shortly after) ten minutes to eight & $7:50 (+\epsilon)$\\\hline
        \bt{Um fast drei viertel acht} & at around \ut{three quarters after eight} & $8:45$
    \end{tabular}
\end{example}
The standard time question is: 
\begin{setdef}
\bt{Wie spät ist es?} : \ut{How late is it?}
\end{setdef}
Calendar units have the following substitutes:
\begin{setdef}
    \begin{tabular}{c|c|c|c}
        \bt{vor} & before & \bt{nach} & after\\\hline
        \multicolumn{4}{c}{Uncertainty markers}\\\hline
        \bt{einige Tage} & Some days & \bt{ungefähr} & approximately
    \end{tabular}
\end{setdef}
The date preposition is \bt{am = an + dem}, since every single month and date is grammatically \bt{männlich}:
\begin{example}
    \begin{tabular}{m{7cm}|m{7cm}}
        \bt{Am zwanzigsten Februar 2019\ldots} & On the twentieth February 2019\ldots\\\hline
        \bt{Der dritte Oktober ist der Tag der deutschen Einheit} & Third October is the day of german unity
    \end{tabular}
\end{example}
The standard date question is:
\begin{setdef}
    \bt{Welches Datum / Welcher Tag ist heute?} : Which date / day is today?
\end{setdef}